\documentclass{beamer}
\usepackage{graphicx}

\usepackage{algorithmicx}
\usepackage{algorithm2e}

\usepackage{amssymb} %Ajoute des symboles mathématiques comme R, Z, forall, exists, etc....
\usepackage{amsmath}
\usepackage{amsthm}

\usepackage{pgf} %permet d'introduire des images vectorielles pour les présentations 
\usepackage{tikz} %permet de créer des graphiques dans le texte 

\usepackage{enumitem} %permet d'énumerer avec un label en (a), (b), plutot que 1. 2. etc... 

\usetheme{Frankfurt}



% Ce projet va servir pour comprendre le Beamer et TikZ tout en révisant les probas de premières années. 
% Plus particulièrement les quelques démos importantes
\renewcommand{\proofname}{Démonstration :} %Commande qui permet de renommer le block "Theorem" par "Démonstrations"


\title{\textbf{Quelques démonstrations importantes de Probabilités de premières années}}
\author{Derrez Lorenzo}

\begin{document}

\begin{frame}
    \maketitle
    \centering
\end{frame}

\begin{frame}{\textbf{Formule des Probabilités totales.}}
    \begin{theorem}
        Soient $(\Omega, \mathbf{P})$ un espace probabilisé fini et $(A_1, \ldots, A_n)$ un système complets d'évènements. Pour tout évènement $B$ :
        \[
            \mathbf{P}(B) = \displaystyle\sum_{i=1}^{n} \mathbf{P}(B \cap A_i) = \displaystyle\sum_{i=1}^{n} \mathbf{P}(B | A_i)\times \mathbf{P}(A_i)
        \]
    \end{theorem}
\end{frame}

\begin{frame}{\textbf{Formule des Probabilités totales.}}
    \begin{proof}
        
    \end{proof}
\end{frame}

\begin{frame}{\textbf{Formule de Bayes.}}
    \begin{theorem}
        Soit $(\Omega, \mathbf{P})$ un espace probabilisé fini.
        \begin{enumerate}[label=(\alph*)]
            \item Si $A$ et $B$ sont deux évènements négligeables tels que $\mathbf{P}(A) > 0 et \mathbf{P}(B) > 0$, alors : 
            
            \[
                \mathbf{P}(A|B) = \mathbf{P}(B|A) \times \frac{\mathbf{P}(A)}{\mathbf{P}(B)}    
            \]
            \item Soient $(A_1, \ldots, A_n)$ un système complet d'évènements et $B$ un évènement tel que $\mathbf{P}(B) > 0$. Alors :
            \[
                \forall i \in [[1, n]] \qquad \mathbf{P}(A_i | B) = \frac{\mathbf{P}(B|A_i)\times \mathbf{P}(A_i)}{\displaystyle\sum_{j=1}^{n}\mathbf{P}(B | A_j)\times \mathbf{P}(A_j)}
            \]
        \end{enumerate}
    \end{theorem}
\end{frame}

\begin{frame}{\textbf{Formule de Bayes.}}
    \begin{proof}
        
    \end{proof}
\end{frame}

\begin{frame}{\textbf{Inégalité de Markov.}}
    \begin{theorem}
        Soit $(\Omega, \mathbf{P})$ un espace probabilisé fini, $X : \Omega \longrightarrow \mathbb{R}$ une variable aléatoire telle que $X \geq 0$ et $a > 0$. Alors : 
        \[
            \mathbf{P}(X \geq a) \leq \frac{\mathbf{E}(X)}{a}.
        \]
    \end{theorem}
\end{frame}

\begin{frame}{\textbf{Inégalité de Markov.}}
    \begin{proof}
        
    \end{proof}
\end{frame}

\begin{frame}{\textbf{Inégalité de Bienaymé-Tchebychev.}}
    \begin{theorem}
        Soient $(\Omega, \mathbf{P})$ un espace probabilisé fini, $X : \Omega \longrightarrow \mathbb{R}$ une variable aléatoire et $\epsilon > 0$. Alors :
        \[
            \mathbf{P}(| X - \mathbf{E}(X) | \geq \epsilon) \leq \frac{\mathbf{V}(X)}{\epsilon^2}.
        \]
    \end{theorem}
\end{frame}

\begin{frame}{\textbf{Inégalité de Bienaymé-Tchebychev.}}
    \begin{proof}
        
    \end{proof}
\end{frame}

\end{document}